%%%%%%%%%%%%%%%%%%%%%%%%%%%%%%%%%%%%%%%%%%%%%%%%%%%%%%%%%%%%%%%%%%%%%%%%%%%
\chapter{Chunking Strategies}
\label{ch:chunking}
%%%%%%%%%%%%%%%%%%%%%%%%%%%%%%%%%%%%%%%%%%%%%%%%%%%%%%%%%%%%%%%%%%%%%%%%%%%

Chunking is the process of splitting documents into smaller pieces before embedding and
indexing. It is one of the highest-leverage decisions in a RAG system---more so than model
choice in many cases. Poor chunking degrades retrieval quality regardless of how good the
embedding model is.

%%%%%%%%%%%%%%%%%%%%%%%%%%%%%%%%%%%%%%%%%%%%%%%%%%%%%%%%%%%%%%%%%%%%%%%%%%%
\section{Why Chunk at All?}
%%%%%%%%%%%%%%%%%%%%%%%%%%%%%%%%%%%%%%%%%%%%%%%%%%%%%%%%%%%%%%%%%%%%%%%%%%%

Two reasons force chunking:

\textbf{1. Token limits.} Embedding models have a maximum input length (typically 512--8{,}192
tokens). A document longer than that limit must be split before it can be embedded. Truncation
is the naive alternative---it silently discards everything past the limit.

\textbf{2. Retrieval precision.} A vector representing a 20{,}000-word document compresses too
much information into one point. When you retrieve that document in response to a specific
question, most of its content is irrelevant to the query. The retrieved context sent to the LLM
is bloated with noise.

Chunking trades recall (a whole document contains the answer somewhere) for precision (the
specific chunk most relevant to the query is returned).

\begin{tipbox}[When Not to Chunk]
Long-context generative models (Gemini 2.0, Claude with 200K context) can ingest entire
documents. If the corpus is small enough that all documents fit in the model's context window,
chunking adds complexity without benefit. Chunking makes sense when: the corpus is too large
to fit in one prompt, retrieval must select relevant content from many documents, or the LLM
context budget is constrained.
\end{tipbox}

%%%%%%%%%%%%%%%%%%%%%%%%%%%%%%%%%%%%%%%%%%%%%%%%%%%%%%%%%%%%%%%%%%%%%%%%%%%
\section{The Chunk Size Tradeoff}
%%%%%%%%%%%%%%%%%%%%%%%%%%%%%%%%%%%%%%%%%%%%%%%%%%%%%%%%%%%%%%%%%%%%%%%%%%%

Chunk size is the central tuning parameter:

\begin{center}
\begin{tabular}{ll}
\toprule
\textbf{Smaller chunks} & \textbf{Larger chunks} \\
\midrule
Higher retrieval precision & Higher recall \\
Less context per retrieved chunk & More context per chunk \\
More chunks to index and store & Fewer chunks, smaller index \\
Better for fact lookup, Q\&A & Better for synthesis, summarization \\
Risk: answer split across boundary & Risk: irrelevant content dilutes chunk \\
\bottomrule
\end{tabular}
\end{center}

Typical defaults by use case:

\begin{center}
\begin{tabular}{ll}
\toprule
\textbf{Use case} & \textbf{Chunk size} \\
\midrule
Q\&A over technical docs & 256--512 tokens \\
General knowledge retrieval & 512--1{,}024 tokens \\
Long-form synthesis & 1{,}024--2{,}048 tokens \\
Code retrieval & Function/class boundary (variable) \\
\bottomrule
\end{tabular}
\end{center}

There is no universal optimal. Test on representative queries from your domain.

%%%%%%%%%%%%%%%%%%%%%%%%%%%%%%%%%%%%%%%%%%%%%%%%%%%%%%%%%%%%%%%%%%%%%%%%%%%
\section{Strategy 1: Fixed-Size Chunking}
%%%%%%%%%%%%%%%%%%%%%%%%%%%%%%%%%%%%%%%%%%%%%%%%%%%%%%%%%%%%%%%%%%%%%%%%%%%

Split every $N$ characters (or tokens), regardless of content boundaries.

\begin{Verbatim}
document = "The model uses attention mechanisms. Attention allows..."

chunk 1 = "The model uses attention mec"   ← breaks mid-word
chunk 2 = "hanisms. Attention allows..."
\end{Verbatim}

\textbf{Advantage:} Simple. Predictable chunk sizes. Easy to implement.

\textbf{Disadvantage:} Breaks at arbitrary points---mid-sentence, mid-word, mid-concept.
Fixed-size chunking without overlap is rarely useful in practice.

%%%%%%%%%%%%%%%%%%%%%%%%%%%%%%%%%%%%%%%%%%%%%%%%%%%%%%%%%%%%%%%%%%%%%%%%%%%
\section{Strategy 2: Fixed-Size with Overlap}
%%%%%%%%%%%%%%%%%%%%%%%%%%%%%%%%%%%%%%%%%%%%%%%%%%%%%%%%%%%%%%%%%%%%%%%%%%%

Like fixed-size, but each chunk shares some content with the previous and next chunk.

\begin{Verbatim}
chunk_size = 500 tokens
overlap    = 100 tokens

chunk 1: tokens 0–499
chunk 2: tokens 400–899    ← shares 100 tokens with chunk 1
chunk 3: tokens 800–1299   ← shares 100 tokens with chunk 2
\end{Verbatim}

\textbf{Why overlap exists:} If the relevant answer spans a chunk boundary, overlap ensures
the answer appears whole in at least one chunk. Without overlap, a sentence split at the
boundary would be partially in chunk~N and partially in chunk~N+1, and neither chunk would
retrieve well.

\textbf{Disadvantage:} Increases index size proportional to overlap fraction. Same content
appears multiple times in search results, requiring deduplication.

\textbf{Typical overlap:} 10--20\% of chunk size. Higher overlap for fine-grained Q\&A;
lower for storage-constrained deployments.

%%%%%%%%%%%%%%%%%%%%%%%%%%%%%%%%%%%%%%%%%%%%%%%%%%%%%%%%%%%%%%%%%%%%%%%%%%%
\section{Strategy 3: Sentence-Boundary Chunking}
%%%%%%%%%%%%%%%%%%%%%%%%%%%%%%%%%%%%%%%%%%%%%%%%%%%%%%%%%%%%%%%%%%%%%%%%%%%

Accumulate text until reaching the size target, but only split at sentence boundaries
(`.`, `!`, `?`, paragraph breaks).

\textbf{Advantage:} Each chunk is semantically coherent. No broken sentences. Better
embedding quality because the model receives grammatically complete input.

\textbf{Disadvantage:} Chunk sizes vary. A single very long sentence can exceed the target.
Requires a sentence tokenizer (spaCy, NLTK, or heuristic regex).

This is the preferred baseline for most RAG systems. memvid approximates this with its
``natural breakpoint'' look-ahead (finding sentence/paragraph endings near the target
boundary).

%%%%%%%%%%%%%%%%%%%%%%%%%%%%%%%%%%%%%%%%%%%%%%%%%%%%%%%%%%%%%%%%%%%%%%%%%%%
\section{Strategy 4: Recursive Character Splitting}
%%%%%%%%%%%%%%%%%%%%%%%%%%%%%%%%%%%%%%%%%%%%%%%%%%%%%%%%%%%%%%%%%%%%%%%%%%%

Split on the most natural boundary available, falling back to coarser splits only when
necessary.

\begin{Verbatim}
Priority order: paragraph → sentence → word → character

1. Try splitting on "\n\n" (paragraph breaks)
2. If any resulting chunk is too large, split on ". " (sentences)
3. If still too large, split on " " (words)
4. If still too large, split on "" (characters)
\end{Verbatim}

This is the default strategy in LangChain and LlamaIndex. It preserves document structure
when possible and only resorts to arbitrary splits when forced.

\textbf{Advantage:} Produces coherent chunks across a wide range of document types without
per-format configuration.

\textbf{Disadvantage:} Can still break logical units (a numbered list, a table row, a code
block) when the target size forces it.

%%%%%%%%%%%%%%%%%%%%%%%%%%%%%%%%%%%%%%%%%%%%%%%%%%%%%%%%%%%%%%%%%%%%%%%%%%%
\section{Strategy 5: Structure-Aware Chunking}
%%%%%%%%%%%%%%%%%%%%%%%%%%%%%%%%%%%%%%%%%%%%%%%%%%%%%%%%%%%%%%%%%%%%%%%%%%%

Use knowledge of the document's format to split on meaningful boundaries:

\begin{itemize}
    \item \textbf{Markdown}: split on headers (\texttt{\#\#}, \texttt{\#\#\#}), preserve
          fenced code blocks and tables whole
    \item \textbf{HTML}: split on \texttt{<section>}, \texttt{<article>}, \texttt{<h2>} tags
    \item \textbf{PDF with structure}: split on section headers from the PDF outline
    \item \textbf{Code}: split on function/class definitions, not arbitrary line counts
    \item \textbf{Tables}: keep each row (or the whole table) as a unit; propagate headers
          to every chunk
\end{itemize}

\textbf{Advantage:} Chunks align with the document's own logical structure.

\textbf{Disadvantage:} Requires a parser per document type. Section-level chunks can be very
long or very short depending on the document.

memvid activates structure-aware mode when it detects Markdown tables or code fences,
propagating table headers to continuation chunks.

%%%%%%%%%%%%%%%%%%%%%%%%%%%%%%%%%%%%%%%%%%%%%%%%%%%%%%%%%%%%%%%%%%%%%%%%%%%
\section{Strategy 6: Semantic Chunking}
%%%%%%%%%%%%%%%%%%%%%%%%%%%%%%%%%%%%%%%%%%%%%%%%%%%%%%%%%%%%%%%%%%%%%%%%%%%

Split based on \textbf{topic shift} rather than size targets. Embed each sentence, compute
cosine similarity between adjacent sentence vectors, and insert a chunk boundary where
similarity drops sharply (a semantic ``cliff'').

\begin{Verbatim}
sentence 1 → vector
sentence 2 → vector  similarity(1,2) = 0.92  → no split
sentence 3 → vector  similarity(2,3) = 0.89  → no split
sentence 4 → vector  similarity(3,4) = 0.41  → SPLIT HERE
sentence 5 → vector  similarity(4,5) = 0.88  → no split
\end{Verbatim}

\textbf{Advantage:} Chunk boundaries align with actual topic transitions. Each chunk has a
single coherent topic. Retrieval precision improves.

\textbf{Disadvantage:} Requires running the embedding model during indexing (not just at query
time). Chunk sizes are unpredictable. Topic shifts within a long sentence are not detected.

Semantic chunking is most valuable for long, multi-topic documents (books, reports, wikis).
For short uniform documents (product descriptions, Q\&A pairs) it adds cost without benefit.

%%%%%%%%%%%%%%%%%%%%%%%%%%%%%%%%%%%%%%%%%%%%%%%%%%%%%%%%%%%%%%%%%%%%%%%%%%%
\section{Strategy 7: Parent-Child (Small-to-Big) Chunking}
%%%%%%%%%%%%%%%%%%%%%%%%%%%%%%%%%%%%%%%%%%%%%%%%%%%%%%%%%%%%%%%%%%%%%%%%%%%

Index small chunks for retrieval precision, but return the larger parent chunk to the LLM
for context richness.

\begin{Verbatim}
Parent chunk:  full section (1200 tokens)  ← not directly indexed
Child chunks:  each paragraph (~200 tokens) ← these are embedded and indexed

Query → retrieve child chunks → return their parent to the LLM
\end{Verbatim}

\textbf{Why this works:} Small child chunks match queries precisely. But the LLM benefits
from the broader context of the parent section, which includes surrounding sentences that
explain the child chunk.

\textbf{Variants:}

\begin{itemize}
    \item Sentence $\to$ paragraph (retrieve sentences, return paragraphs)
    \item Paragraph $\to$ section (retrieve paragraphs, return sections)
    \item Chunk $\to$ full document (retrieve chunks, return whole document)
\end{itemize}

memvid implements this: \texttt{FrameRole::DocumentChunk} for child chunks,
\texttt{FrameRole::Document} for the parent, linked by \texttt{parent\_id}.

\textbf{Disadvantage:} Doubles the storage (parent and child both stored). The parent context
may still include irrelevant content. Works best when document sections have natural boundaries.

%%%%%%%%%%%%%%%%%%%%%%%%%%%%%%%%%%%%%%%%%%%%%%%%%%%%%%%%%%%%%%%%%%%%%%%%%%%
\section{Strategy 8: Late Chunking}
%%%%%%%%%%%%%%%%%%%%%%%%%%%%%%%%%%%%%%%%%%%%%%%%%%%%%%%%%%%%%%%%%%%%%%%%%%%

An emerging strategy for ColBERT-style models. Rather than chunking before embedding, embed
the whole document and then extract chunk-level embeddings from the token-level outputs.

\begin{Verbatim}
Document → [embedding model] → token vectors (one per token)
                              → pool each chunk's token range into one vector
\end{Verbatim}

\textbf{Advantage:} Eliminates boundary artifacts---a sentence at the start of the document
has the same embedding quality as one in the middle.

\textbf{Disadvantage:} Requires a model that exposes token-level outputs (ColBERT-style or raw
transformer outputs). Not supported by the standard OpenAI-compatible embeddings API.
Currently experimental.

%%%%%%%%%%%%%%%%%%%%%%%%%%%%%%%%%%%%%%%%%%%%%%%%%%%%%%%%%%%%%%%%%%%%%%%%%%%
\section{How memvid Handles Chunking}
%%%%%%%%%%%%%%%%%%%%%%%%%%%%%%%%%%%%%%%%%%%%%%%%%%%%%%%%%%%%%%%%%%%%%%%%%%%

memvid's chunking implementation (\texttt{src/memvid/chunks.rs}):

\begin{itemize}
    \item Default: 1{,}200-character chunks with $\pm$240-char look-ahead for natural breakpoints
    \item Documents under 2{,}400 chars are not chunked
    \item Structure-aware mode activates on detection of Markdown tables or code fences
    \item Tables: splits between rows, propagates headers to continuation chunks
    \item Code: keeps fenced blocks whole where possible
    \item Parent-child: \texttt{FrameRole::Document} links to \texttt{FrameRole::DocumentChunk}
          children via \texttt{parent\_id}; \texttt{TextChunkManifest} records each child's
          \texttt{(start, end)} character range
\end{itemize}

\textbf{Gaps:} no overlap, no semantic chunking, no sentence-boundary splitting (heuristic
only), no recursive character fallback. These are reasonable defaults for a general-purpose
embedded store but may need supplementing for domain-specific corpora.

%%%%%%%%%%%%%%%%%%%%%%%%%%%%%%%%%%%%%%%%%%%%%%%%%%%%%%%%%%%%%%%%%%%%%%%%%%%
\section{Practical Checklist}
%%%%%%%%%%%%%%%%%%%%%%%%%%%%%%%%%%%%%%%%%%%%%%%%%%%%%%%%%%%%%%%%%%%%%%%%%%%

When designing a chunking strategy:

\begin{enumerate}
    \item \textbf{Know your documents.} Short uniform records behave differently from long
          heterogeneous documents. Sample 20--30 representative documents before choosing
          a strategy.

    \item \textbf{Know your queries.} Fact-lookup queries benefit from small precise chunks.
          Synthesis queries benefit from larger chunks.

    \item \textbf{Start simple.} Sentence-boundary chunks with 15\% overlap is a good default.
          Measure retrieval quality before adding complexity.

    \item \textbf{Measure chunk size distribution.} Check the p50, p95, and max chunk sizes.
          Outliers (single-sentence chunks, oversized sections) degrade retrieval differently
          than average chunks.

    \item \textbf{Always store metadata.} At minimum: source, section, date. This is cheap to
          add and expensive to reconstruct later.

    \item \textbf{Plan for re-chunking.} When you change chunking strategy or embedding model,
          the entire index must be rebuilt. Design ingestion pipelines so re-chunking is not a
          one-time manual operation.
\end{enumerate}
